\documentclass{dsekprotokoll}

\usepackage[T1]{fontenc}
\usepackage[utf8]{inputenc}
\usepackage[swedish]{babel}
\usepackage{multicol}
\usepackage{hyperref}

\hypersetup{
    colorlinks=true,
    linkcolor=blue,
    urlcolor=blue
}

\setheader{Riktlinje för Arkivering}{Riktlinje}{\today}

\title{Riktlinje för Arkivering}
\author{}

\begin{document}
\maketitle

\section{Formalia}
\subsection{Sammanfattning}
Denna riktlinje ger direktiv kring hur material i D-sektionens verksamhet ska arkiveras.

\subsection{Syfte}
Syftet med denna riktlinje är att säkerställa att material av framtida intresse för sektionen bevaras och att det organiseras på ett systematiskt och lättillgängligt vis.

\subsection{Omfattning}
Riktlinjen omfattar hela D-sektionens organisation.

\subsection{Historik}
Utkast färdigställt av
\begin{itemize}
    \item Arkivarie 2015--2023 Rasmus Göransson
    \item Arkivarier 2023--2024 Jacob Annefors, Philip Nielsen och Mikolaj Sinicka
    \item Arkivarier 2024 Axel Hartwig och Jacob Johansson
    \item Vice Informationsansvarig 2024 Anahita Chavan
    \item Informationsansvarig 2024 Dag Hemberg
\end{itemize}

Riktlinje antagen på S21 2024

\section{Grundläggande arkivering}
I arkivet ska material som producerats vid D-sektionen och som kan anses vara till gagn för eftervärlden samlas, liksom sådant som överlämnats eller inkommit till D-sektionen och som inte arkiveras någon annanstans. Inom D-sektionen ska producerat material överlämnas kontinuerligt till arkivet. Handlingar, protokoll och dylikt för sektionsmötena, studierådet, styrelsen, och liknande överlämnas företrädesvis i sorterat skick till arkivet i slutet av verksamhetsåret.

Vid överlämning av digitalt material bör detta i första hand skickas in till Arkivarierna via arkiv@dsek.se, och i andra hand till Informationsansvarig via infa@dsek.se. Vid överlämning av analogt material bör detta lämnas direkt till arkivarierna vid ett lämpligt tillfälle. För övrigt material bör arkivarier kontaktas i förväg för att komma överens om lämpligt tillvägagångssätt för inlämning av materialet. 

% how to send in
% Analogt material är exempelvis: Bilder, Affischer, Kläder, Märken, m.m. Dessa läggs i brevlådan i mån av plats eller lämnas direkt till arkivarier.

\section{Definitioner}

\subsection{D-sektionens producerade material}
Med D-sektionens producerade material menas nedan beskrivet digitalt och analogt material som producerats till D-sektionens event samt övrigt material som producerats av D-sektionens utskott, underutskott och projektgrupper.

\subsubsection*{Digitalt material}
Med digitalt material menas: 
\begin{itemize}
    \item Bilder och videoklipp
    \item Filmer
    \item Media från sektionens officiella sociala medie-konton
    \item Digitala versioner av affischer, sektionstidningar och dokument
\end{itemize}
% Dessa skickas in till \href{mailto:arkivarie@dsek.se}{arkivarie@dsek.se}.

\subsubsection*{Analogt material}
Med analogt material menas: 
\begin{itemize}
    \item Märken
    \item Allt tryckt material
    \item Klädesplagg med undantag för funktionärsspecifika kläder
    \item Merch
    \item Viktiga dokument som inte är möteshandlingar
\end{itemize}

\subsubsection*{Övrigt material}
Med övrigt material menas exempelvis:
\begin{itemize}
    \item Stulna klenoder från andra sektioner
    \item Jubileumsmaterial
    \item Krimskrams
\end{itemize}

\subsection{Möteshandlingar}
Endast protokoll för sektionsmöten och styrelsemöten ska arkiveras analogt enligt nedan, övriga dokument arkiveras digitalt.

\subsubsection*{Sektionsmöte}
\begin{itemize}
    \item Kallelser
    \item Protokoll
    \item Handlingar
\end{itemize}

\subsubsection*{Styrelsemöte}
\begin{itemize}
    \item Kallelser
    \item Protokoll
    \item Handlingar
    \item Per Capsulam-beslut
\end{itemize}

\subsubsection*{Studierådsmöte}
\begin{itemize}
    \item Kallelser
    \item Protokoll
    \item Handlingar
\end{itemize}

\subsubsection*{Övriga Möteshandlingar}

Om sektionen fattar beslut eller genomför åtgärder som arkivet bedömer vara av långsiktigt värde, bör dessa bevaras för framtiden. Exempelvis innefattar dessa viktiga beslut inom husstyrelsen eller större möten med E-styrelsen.

\subsection{Inkommet material}
Inkommet material omfattar till D-sektionen överlämnade presenter, skrivelser ställda till D-sektionen, inbjudningar och annat. Under denna rubrik räknas även artiklar, annonser och notiser rörande D-sektionen.

\subsection{Ekonomi och bokföring}
Arkivering av sektionens ekonomi och bokföring sköts av Skattmästeriet i enlighet med sektionens övriga styrdokument.

\subsection{Övrigt material}
Vidare sparas övrigt material som kan vara av intresse för framtiden enligt arkivets bedömning. Hit hör bland annat material av fånig och studentikos karaktär som livar upp arkivariernas arbete.

\section{Övrigt}

\subsection*{Kvantiteter}

Överlämning av föremål i den nedan följande listan bör innefatta minst 3 exemplar.

\begin{itemize}
    \item Märken
    \item Tryckt material
    \item Merchartiklar
    \item Klädesplagg
\end{itemize}

\end{document}
